% !TEX root = ./manuscript.tex
\documentclass[pdflatex,sn-mathphys-num]{sn-jnl}

\usepackage{amsmath,amssymb}
\usepackage{booktabs}
\usepackage{array}
\usepackage{xurl}

\newcolumntype{Z}[1]{>{\raggedright\arraybackslash}p{#1}}

\raggedbottom

\begin{document}

\title[Structural Labels for Stratified Audit Budget Allocation in Knowledge-Graph Dependency Flows]{Structural Labels for Stratified Audit Budget Allocation in Knowledge-Graph Dependency Flows}

\author*[1]{\fnm{Haoran} \sur{Ma}}\email{mahaoran0000@foxmail.com}
\author[1,2]{\fnm{Ningning} \sur{Wang}}\email{wangningning@bistu.edu.cn}

\affil*[1]{\orgdiv{College of Management Science and Engineering}, \orgname{Beijing Information Science and Technology University}, \orgaddress{\city{Beijing}, \postcode{102200}, \country{China}}}
\affil[2]{\orgdiv{Institute of Information Systems, ESG Intelligent Application Innovation Research Center}, \orgname{Beijing Information Science and Technology University}, \orgaddress{\city{Beijing}, \postcode{102200}, \country{China}}}

\abstract{Knowledge-graph and retrieval-augmented information systems expose dependency flows, entity bindings, and verification records that must be audited under fixed budgets. Prior score-centered formulations make this task look like another ranking problem, but a scalar score does not tell curators which artifacts are bottlenecks, which evidence is redundant, or how downstream records are affected. We recast the framework around structural labels rather than score improvement. The first contribution is a structural label extractor, not a scorer, that converts graph topology into audit-ready boolean flags for bottleneck and redundancy. The second contribution is a Life-Saving First stratified budget allocation policy guided by these labels and evaluated with Impact Coverage@K over reachable descendants in KG dependency flows. On the Countries-KG semantic fixture, the extractor recovers controlled bottleneck and redundancy labels while an undirected Semantic-Similarity Baseline (TF-IDF) cannot recover the directed-flow labels. In the cached-label audit case, the stratified policy covers transitive downstream impact under the same budget and is compared against Flat Top-K using the shared raw risk score. The evidence is deliberately bounded: it is a clean semantic fixture and seeded scalability test, not production knowledge-base validation, human usefulness evidence, arbitrary-noise robustness, or causal identification.}

\keywords{knowledge graphs, structural labels, audit budget allocation, impact coverage, information systems, dependency flows}

\maketitle

\section{Introduction}\label{sec:introduction}

Knowledge-intensive information systems increasingly expose intermediate records rather than only final answers. A knowledge-graph workflow may expose entity nodes, relation paths, retrieval passages, verification steps, and dependency links. These artifacts are useful for maintenance only when a curator can decide what to inspect under a fixed budget. The difficulty is not simply that some artifacts receive low scalar scores. A curator also needs to know whether an artifact is a bottleneck for later records, whether another artifact provides redundant coverage, and which downstream nodes would be affected if the selected artifact is audited.

This paper therefore moves the contribution away from a scoring race. We propose a \emph{structural label extractor} that turns directed graph topology into audit-ready boolean fields: \texttt{bottleneck} and \texttt{redundancy}. These fields are not reward-model predictions and are not claimed to improve task reasoning. They are reusable labels for audit records in dependency-flow settings.

The second contribution is a \emph{stratified budget allocation} policy, Life-Saving First. The policy first allocates budget to Critical Bottleneck nodes, then to Unique Evidence nodes, then to Redundancy Group Samples, and finally to Fallback nodes. Raw risk score is used exclusively as a tie-breaker within the same stratification layer, not as the primary allocation driver. Flat Top-K uses the shared \detokenize{raw_risk_score} as the sole ranking criterion, where the KG setting defines it as trace-local min--max normalized downstream impact count.

The evaluation follows this reframing. Table~\ref{tab:label-validation} reports Countries-KG structural label extraction, with Countries-KG used as a clean semantic fixture containing 30 entities and 189 triples. Table~\ref{tab:impact-coverage} reports Impact Coverage@K rather than Recall@25-style ranking. Impact Coverage counts all reachable descendants (transitive closure) from selected nodes, not only direct successors. This reflects the maintenance reality that a root bottleneck can affect a long downstream chain, and auditors need visibility into the entire affected subgraph.

\section{Method}\label{sec:method}

\subsection{Structural label extractor}

Given a directed dependency-flow graph $G=(V,E)$, the extractor emits two boolean labels per node. A node is marked as a bottleneck when its removal lowers reachable terminal-flow coverage from the frozen source set or when it occupies a critical source-to-terminal dependency position. A node is marked as redundant when its dependency coverage overlaps another node by Jaccard similarity greater than $\theta=0.85$. We define redundancy gold labels using Jaccard $>0.85$ over dependency coverage, acknowledging that strict identical coverage is rarely observed in semantic KGs.

The TF-IDF comparator is named Semantic-Similarity Baseline (TF-IDF), not a graph-based method. This baseline is included to show that undirected semantic similarity, which lacks flow direction, cannot recover structural bottleneck/redundancy labels.

\subsection{Life-Saving First budget allocation}

Budget allocation proceeds sequentially. If a layer cumulatively exceeds K, selection stops within that layer sorted by impact; subsequent layers are not invoked for this trace. Critical Bottleneck nodes can override the ordinary auditable-node filter because root or scaffold bottlenecks may affect the whole downstream subgraph. Ordinary baselines rank only the auditable node pool.

The four strata are: Critical Bottleneck, Unique Evidence, Redundancy Group Samples, and Fallback. Within a stratum, nodes are ordered by downstream impact and then by the shared raw risk score only to break ties. The No-Fallback Ablation uses the same strata but randomizes same-layer tie breaking, testing whether Impact Coverage depends on the scalar fallback.

\section{Results}\label{sec:results}

\begin{table}[t]
\caption{Countries-KG structural label F1 comparison. The semantic fixture contains 30 entities and 189 triples; labels are cached with seed 20260711. The TF-IDF row is the Semantic-Similarity Baseline (TF-IDF). This baseline is included to show that undirected semantic similarity, which lacks flow direction, cannot recover structural bottleneck/redundancy labels.}
\label{tab:label-validation}
\centering
\small
\begin{tabular}{Z{0.34\linewidth}rrr}
\toprule
\textbf{Method} & \textbf{Bottleneck F1} & \textbf{Redundancy F1} & \textbf{Redundancy positives} \\
\midrule
Structural label extractor & 1.000 & 1.000 & 124 \\
Semantic-Similarity Baseline (TF-IDF) & 0.242 & 0.405 & 124 \\
\bottomrule
\end{tabular}
\end{table}

Table~\ref{tab:label-validation} shows the controlled label-validation result. The structural extractor recovers the graph-defined labels, whereas the undirected semantic baseline fails to recover directed-flow bottleneck and redundancy roles. The result supports label extraction on a clean semantic fixture; it is not a claim about noisy open-domain KGs.

\begin{table}[t]
\caption{Impact Coverage@K of Life-Saving First stratified policy vs. flat Top-K baseline (using the shared \detokenize{raw_risk_score} as the sole ranking criterion), centrality, position, random, and no-fallback ablation.}
\label{tab:impact-coverage}
\centering
\small
\begin{tabular}{Z{0.31\linewidth}rrrr}
\toprule
\textbf{Method} & \textbf{Impact Coverage@K} & \textbf{Avg. path length} & \textbf{Early truncation} & \textbf{Budget used} \\
\midrule
Life-Saving First & 1.000 & 1.000 & 0.000 & 1.000 \\
Flat Top-K & 0.733 & 0.828 & 0.000 & 1.000 \\
Centrality & 0.650 & 0.761 & 0.000 & 1.000 \\
Position & 0.000 & 0.000 & 0.000 & 1.000 \\
Random & 0.333 & 0.622 & 0.000 & 1.000 \\
No-Fallback Ablation & 1.000 & 1.000 & 0.000 & 1.000 \\
\bottomrule
\end{tabular}
\end{table}

Table~\ref{tab:impact-coverage} evaluates budget allocation on cached Countries-KG labels. The comparison is deliberately between stratification and a single-score budget rule that shares the same scalar source. The no-fallback ablation matches the main policy in this fixture, indicating that the observed coverage is driven by the structural strata rather than by scalar fallback ordering. The low early-truncation rate confirms that the policy does not collapse entirely inside the first layer; in this clean fixture, selection primarily uses the critical-bottleneck layer and the redundancy-group-sample layer.

\section{Boundary diagnosis}\label{sec:boundary}

Appendix B reports a Necessary Condition Diagnosis on PRM800K. In the locked process-annotation route, the supervised fidelity field \texttt{w\_struct} reaches Spearman $\rho=0.611$, whereas undirected TF-IDF graph-only necessity collapses to $\rho=0.000$ in the high trace-length proxy. This is not a hidden main result. It defines a boundary condition: structural labels require at least moderately informative directed dependency edges and should not be expected to work when the graph constructor supplies sparse or directionless similarity.

\section{Limitations}\label{sec:limitations}

The semantic main experiment is deliberately bounded to Countries-KG to provide a ground-truth dependency flow for controlled F1 evaluation. We do not claim that structural labels are robust to arbitrary KG noise; rather, we treat this as a proof-of-concept on a clean semantic fixture. Scaling to noisy, open-domain KGs (e.g., via injected edge perturbations) is a necessary future step, not a claim of the current framework.

The study also does not report human usefulness evidence. The submission package contains a blank future human-evaluation protocol, but no returned ratings, evaluator provenance, declarations, timestamps, hashes, or reliability analysis are available. The present evidence therefore concerns structural label extraction and Impact Coverage on KG dependency flows, not human decision time or human audit accuracy.

\section{Conclusion}\label{sec:conclusion}

This paper reframes the SC-FMA line as structural label extraction and label-guided budget allocation. The contribution is not a new universal scorer. It is a graph-topology interface that emits bottleneck and redundancy flags, plus a Life-Saving First policy that uses those flags for fixed-budget audit visibility. The evidence supports controlled label extraction on Countries-KG, seeded scalability in Appendix A, Impact Coverage comparison on cached KG flows, and PRM800K necessary-condition boundaries in Appendix B.

\bibliography{references}
\end{document}
